\section{Negative results and interpretation boundaries}
\label{app:negative-results}

\subsection{Analytic acceleration is a scoped non-result}

% CLAIM: C04
On all three analytic references, PV truth and PVA truth have the same
\AnalyticPVLagMS{}-ms best grid lag, and their position RMSE values differ by no
more than \(\AnalyticPVPVAMaxRMSEDifference\,\si{\radian}\).  This is useful negative
evidence because the target-component comparison is controlled: requested
position and velocity are identical, while requested acceleration differs.
Under these runs, both conditions reach the same requested position and
velocity endpoint, so the acceleration difference does not alter the recorded
position endpoint.  The experiment does not measure internal effort, profile
shape, acceleration-estimator quality, or acceleration-active regimes.
Accordingly it supports ``no observed position improvement here,'' not a
general equivalence or uselessness claim.

\subsection{Development-trace finite differences}

% CLAIM: C06,C07,N01
The complete CSV table in \cref{tab:csv-tracking} preserves causal and
noncausal rows rather than reporting only the best derivative variant.  The
negative result is consistent across the tested matrix: every unfiltered
finite-difference-derived pipeline has higher position RMSE than P under the
fixed projection policy.  Yet the PVA rows and PV rows are not a clean
acceleration-only contrast.  P and PV project no targets, whereas each PVA row
projects \CSVProjectionRate\%; the raw PVA acceleration estimates reach
\SI{\CSVRawAccelerationPeak}{\radian\per\second\squared}.  Projection and
estimation therefore change together.

The scope is intentionally narrow.  The input is one development trace, its
recorded timestamps are ignored, and its rows are treated as uniformly spaced
by \SI{\ControlPeriodMS}{\milli\second}.  There is no derivative truth,
independent trajectory, hardware execution, or irregular-sampling test.
Filtered differentiation, model-based estimation, timestamp-aware processing,
and matched projection ablations remain untested by this result.

\subsection{One-factor limit sensitivity}

The Phase~A limit sweep changes acceleration or jerk one factor at a time.
On the CSV P baseline, reducing the jerk limit from the configured
\SI{\JerkLimit}{\radian\per\second\cubed} to
\SI{\CSVLowJerkSweep}{\radian\per\second\cubed} increases RMSE by a factor of
\CSVLowJerkRMSEFactor{} and adds
\SI{\CSVLowJerkLagIncreaseMS}{\milli\second} of best grid lag.  Raising the jerk limit from the
configured value to
\SI{\CSVHighJerkSweep}{\radian\per\second\cubed} changes the RMSE ratio
to \CSVHighJerkRMSEFactor{} with unchanged best lag.  Because the design is one-factor-at-a-time,
these points neither identify an acceleration--jerk interaction nor recommend
limits outside the configured operating point.
